\scanpage{116}
\mspage{- 121 -}\hfill IX

10. \uline{Propriétés de l'accouplement de monodromie}. \uline{Théorème d'intégrité et de positivité}\unsure{.}

Dans le présent paragraphe, nous allons donner tout d'abord certaines propriétés formelles essentiellement triviales des accouplements de monodromie, le détail (un peu fastidieux) des vérifications étant laissé au lecteur. D'autre part, nous indiquons un théorème plus profond (10.4) concernant ces accouplements, et indiquons comment on peut se réduire pour sa démonstration à un cas spécial de jacobiennes, qui sera traité (par voie transcendante) au par. 12 (qui est d'ailleurs logiquement indépendant du par. 11).

10.1. \uline{Transport de structure}. Il est clair, sur la définition donnée de l'accouplement de monodromie (9.1.2), que cet accouplement est compatible avec le transport de structure en la situation $(S,A_K,A'_K)$.

10.2. \uline{Symétrie}. Si nous partons du schéma dual $A'_K$ de $A_K$ et lui appliquons la construction du par. 9, nous devons introduire le schéma dual $(A'_K)'$ de $A_K$, qui est canoniquement isomorphe à $A_K$ par le théorème de bidualité ([5] ou VIII 3.2). Désignons par $u_\ell^{A_K}$ l'accouplement de monodromie (9.1.2) relativement à $A_K$, de sorte qu'on trouve de même un accouplement
\[
 u_\ell^{A'_K}:\underline M'_\ell\otimes\underline M_\ell\longrightarrow\mathbb Z_\ell\quad.\tag*{(10.2.1)}
\]
Ceci posé, je dis que ce dernier se déduit de $u_\ell^{A_K}$ par la symétrie
\[
 s:\underline M_\ell\otimes\underline M'_\ell\xleftarrow{\sim}\underline M'_\ell\otimes\underline M_\ell\quad,
\]
i.e. on a la formule

\mspage{\handwritten{439}}
% stamp: numbered label 439
% archive sheet number: 115
